\documentclass[12pt]{article}
\usepackage[utf8]{inputenc}
\usepackage[T1]{fontenc}
\usepackage{hyperref}
\title{Wabo Language Resources for K-12 Curriculum Development}
\author{Compiled by Claude Sonnet 5 (claude-sonnet-5) for Andr\'as Kornai}
\date{}
\begin{document}
\maketitle
\section{Introduction}
Wabo (ISO 639-3 wbb) is an Austronesian language of Yapen Island, Papua, Indonesia. Wikipedia gives about 1,500 speakers (1987), and Joshua Project gives 2,700. A 2025 study in Wabompi Village found it definitely endangered by UNESCO indicators and moribund by SIL indicators. Speakers mostly use Indonesian, and in some villages Ambai.
\section{Primary Educational Resources}
No primers, textbooks or school materials in Wabo were found in the searches.
\section{Linguistic and Reference Materials}
\begin{itemize}
\item Wikipedia overview: \url{https://en.wikipedia.org/wiki/Wabo_language}
\item Vitality study, ``Wabo Language in Wabompi Village: Evaluation of Language Vitality and Endangerment'': \url{https://tpls.academypublication.com/index.php/tpls/article/view/12523}
\item News report (June 2026) on a BRIN verification workshop and grammar work, with data archived in the National Scientific Repository and in ELAR: \url{https://news.republika.co.id/berita/tggmej416/}
\end{itemize}
\section{Religious and Community Resources}
\begin{itemize}
\item Joshua Project people profile: \url{https://joshuaproject.net/people_groups/15811}
\end{itemize}
\section{Digital and Audio Resources}
Recordings are reported to be archived in the Endangered Languages Archive (ELAR). I did not see a direct catalogue URL in the search results.
\section{Curriculum Development Context}
The language is in steep decline, with active use reported mainly among grandparents. The current grammar and archival work by BRIN could support later revitalisation materials. No school materials exist yet in the sources found.
\end{document}